\chapter{Problem Definition and Benchmark}
\label{chap:problem}

The planning problem consists of fixed benchmark data, a selected timetable, and
a vehicle schedule that operates exactly the selected trips. The main objects
are the public-transport network, candidate trips, time-window rules,
vehicle-block compatibility, electric-bus autonomy, charging resources, and cost
coefficients.

\section{Study Scope and Input Assumptions}

The solver works on fixed MINOA Senior input files \cite{minoa-io,minoa-problem,minoa-rules}. It does not change the
network, candidate trips, deadhead arcs, vehicle data, charging rates, node
capacities, or cost coefficients. The decisions start after the data are fixed:
the solver selects a feasible timetable from the candidate trips and then
constructs a feasible vehicle schedule for that timetable.

Let \(I\) be one benchmark instance. Its input data can be written as

\begin{equation}
I=
\left(
T,D,N,A^{dh},V,V_E,\mathcal{W},\mathcal{R},\mathcal{K}
\right),
\label{eq:instance-tuple}
\end{equation}

where \(T\) is the candidate trip set, \(D\) is the set of line directions,
\(N\) is the set of nodes, \(A^{dh}\) is the deadhead-arc data, \(V\) is the
set of vehicle types, and \(V_E\subseteq V\) is the set of electric vehicle
types. The tested Senior instances contain one electric type and one ICE type.
The formulation also allows several vehicle types, as permitted by the problem
statement. The set \(\mathcal{W}\) contains time-window dependent rules,
\(\mathcal{R}\) contains charging and resource-capacity data, and
\(\mathcal{K}\) contains the cost coefficients \cite{minoa-problem,minoa-io}.

The output contains two coupled objects:

\begin{equation}
S=(\pi,\Omega),
\label{eq:solution-pair}
\end{equation}

where \(\pi\subseteq T\) is the selected timetable and \(\Omega\) is the vehicle
schedule. A solution is feasible only if the selected timetable and the vehicle
schedule describe the same operated passenger trips:

\begin{equation}
T(\Omega)=\pi .
\label{eq:tt-vs-equality}
\end{equation}

\Cref{eq:tt-vs-equality} is the central link between timetabling and vehicle
scheduling. It means that every selected trip must appear in the vehicle
schedule and no unselected candidate trip may appear there. This prevents the
solver from choosing one timetable and operating a different one. The formal
coverage constraints are stated later in
\Cref{eq:trip-covered-once,eq:unselected-not-covered}.

\section{Network and Service Representation}

\subsection{Lines, Directions, Trips, and Terminals}

The candidate passenger trips are grouped by line direction. For each direction
\(d\in D\), let \(T_d\subseteq T\) be the trips of that direction. The
direction sets form a partition of the candidate pool:

\begin{equation}
T=\bigcup_{d\in D}T_d,
\qquad
T_d\cap T_{d'}=\emptyset
\quad \forall d\neq d'.
\label{eq:trip-partition}
\end{equation}

Each trip \(i\in T\) has a start node \(sn(i)\), an end node \(en(i)\), a start
time \(st(i)\), an end time \(et(i)\), and a main-stop time \(a(i)\). The
headway rules use \(a(i)\) because the problem defines service regularity at a
main stop, not only at the terminal departure time.

The duration of a passenger trip is

\begin{equation}
\tau_i^{trip}=et(i)-st(i),
\qquad
\tau_i^{trip}>0.
\label{eq:trip-duration}
\end{equation}

The planning model also uses the start and end nodes to decide whether two trips
can be connected directly or must be connected through the depot.

\subsection{Depot, Nodes, and Deadhead Movements}

The depot is the location where every vehicle block starts and ends. Terminal
nodes are the locations where passenger trips start or end and where vehicles
may wait, park, or charge. A deadhead arc describes a non-passenger movement.
For a deadhead movement \(a\in A^{dh}\), I write travel time and distance as

\begin{equation}
\tau^{dh}_{a}(t),\qquad \ell^{dh}_{a},
\label{eq:deadhead-time-distance}
\end{equation}

where \(t\) is the relevant time at which the deadhead time window is chosen.
The travel time \(\tau^{dh}_a(t)\) is measured in seconds and may depend on the
time window. The distance \(\ell^{dh}_a\) is measured in kilometres and does
not depend on the time of day \cite{minoa-io}.
Pull-out and pull-in movements are special deadheads:

\begin{equation}
t_b^{po}=\tau^{dh}_{dep,sn(i_1)}(st(i_1)),
\qquad
t_b^{pi}=\tau^{dh}_{en(i_q),dep}(et(i_q)),
\label{eq:pull-times}
\end{equation}

for a block \(b=(i_1,\ldots,i_q)\). The notation \(t_b^{po}\) denotes
pull-out duration from the depot to the first trip, and \(t_b^{pi}\) denotes
pull-in duration from the last trip back to the depot.

\subsection{Planning Horizon and Time Windows}

All activities of a vehicle block must lie inside the planning horizon of the
instance. The horizon is divided into time windows
\(\mathcal{H}=\{1,\ldots,k\}\) by time instants
\(t_0<t_1<\cdots<t_k\). The time-window index is

\begin{equation}
\omega(t)=h
\quad\Longleftrightarrow\quad
t_{h-1}<t\le t_h
\label{eq:time-window-map}
\end{equation}

for the time-window index containing time \(t\). At the first boundary,
\(\omega(t_0)=1\) is used by the implementation, and times outside the horizon
are mapped to the nearest boundary window, matching the input/output
convention for external times \cite{minoa-io}. If \(H=[t^{min},t^{max}]\), then every
activity interval \([s_a,e_a]\) in the vehicle schedule must satisfy

\begin{equation}
t^{min}\le s_a < e_a \le t^{max}.
\label{eq:horizon-activity}
\end{equation}

The headway rules, pull-in times, pull-out times, and stopping-time bounds are
indexed by \(\omega(\cdot)\). For example, the maximum headway of direction
\(d\) in window \(h\) is denoted by \(h^{max}_{d,h}\), while the minimum and
maximum stopping times at node \(n\) are
\(\delta^{min}_{n,h}\) and \(\delta^{max}_{n,h}\). The relevant window for an
in-line stop is selected by the arrival time at the node, here \(et(i)\). The
same time-window choices appear in the in-line and depot-bridge constraints
in \Cref{eq:inline,eq:depot}.

\section{Timetable Selection Layer}

\subsection{Candidate Trips and Selected Timetable}

For each direction \(d\), the selected timetable is a subset
\(\pi_d\subseteq T_d\). The complete selected timetable is

\begin{equation}
\pi=\bigcup_{d\in D}\pi_d .
\label{eq:problem-selected-timetable}
\end{equation}

The problem is not to operate all candidate trips. The input contains many
possible departures, and the solver selects enough of them to satisfy the
service rules. The selected-trip percentage is therefore

\begin{equation}
p_T=\frac{|\pi|}{|T|}\cdot 100\%.
\label{eq:problem-selected-share}
\end{equation}

The selected-trip percentage measures the density of the final timetable
relative to the candidate pool.

\subsection{Headway Feasibility}
\label{sec:headway-feasibility}

For two consecutive selected trips \(i,j\in\pi_d\), the headway is measured at
the main stop:

\begin{equation}
w_{ij}=a(j)-a(i).
\label{eq:problem-headway}
\end{equation}

If \(I^{max}_{ij}\) is the maximum permitted headway for this pair and time
window, the feasibility condition is

\begin{equation}
0<w_{ij}\le I^{max}_{ij}.
\label{eq:problem-headway-bound}
\end{equation}

The value \(I^{max}_{ij}\) follows the border-effect rule in the MINOA problem
statement. If both main-stop times lie in the same time window, the maximum
headway of that window is used. If the two trips lie in two consecutive time
windows, the allowed bound is the larger of the two adjacent maximum headways.
This avoids making a pair infeasible only because the two main-stop times fall
on different sides of a time-window border. The same rule is used in the formal
headway graph in \Cref{eq:headway-edge,eq:headway-bound}.

Equivalently, each direction can be represented by a directed headway graph

\begin{equation}
H_d=(T_d,E_d),
\qquad
(i,j)\in E_d
\Longleftrightarrow
a(i)<a(j),\; a(j)-a(i)\le I^{max}_{ij}.
\label{eq:problem-headway-graph}
\end{equation}

The selected timetable \(\pi_d\) is then an initial-to-final path in
\(H_d\). This graph interpretation is important because it connects the
passenger-service rule to the algorithmic timetable-selection step. Graph and
network representations are also standard tools in vehicle scheduling and
network-flow models \cite{kliewer2006,ford1962,lawler1976}.

\begin{figure}[H]
\centering
  \includegraphics[width=\textwidth]{figures/fig36_headway_selection.png}
\caption{Timetable selection under headway constraints.}
\label{fig:headway-selection}
\end{figure}
\FloatBarrier

\Cref{fig:headway-selection} shows the idea for one direction. Grey departures
are available alternatives, while the blue departures form the selected
timetable path. Consecutive blue trips must stay within the maximum headway,
and the first and last selected trips must satisfy the boundary requirements.
The figure explains why timetable selection is a graph path problem in this
thesis: the chosen blue trips form a feasible path through the candidate
departures, and every later vehicle block is built only from this selected path.
This is the first decision layer of the problem: before any vehicle can be
scheduled, the solver must choose a passenger timetable that does not leave
excessive gaps between consecutive departures in one direction.

\subsection{Initial and Final Trip Requirements}

The selected timetable must also satisfy boundary requirements. Let
\(T_d^{start}\subseteq T_d\) be the admissible initial trips and
\(T_d^{end}\subseteq T_d\) be the admissible final trips for direction \(d\).
If \(\pi_d=(i_1,\ldots,i_q)\), then

\begin{equation}
i_1\in T_d^{start},
\qquad
i_q\in T_d^{end}.
\label{eq:boundary-trips}
\end{equation}

Thus, timetable selection is not only a sparse subset choice. It is a path
choice:

\begin{equation}
\pi_d=(i_1,\ldots,i_q)
\quad\text{with}\quad
(i_r,i_{r+1})\in E_d
\;\; r=1,\ldots,q-1.
\label{eq:timetable-path-definition}
\end{equation}

\Cref{chap:model} uses this path view in the formal selected-timetable
definition in \Cref{eq:selected-trips} and in the dynamic recursion
\Cref{eq:timetable-dp}.

\section{Vehicle-Scheduling Layer}

\subsection{Vehicle Blocks and Trip Coverage}

A vehicle block is an ordered sequence of vehicle activities. In its passenger
trip view, a block can be written as

\begin{equation}
b=(i_1,i_2,\ldots,i_q),
\qquad i_r\in\pi.
\label{eq:block-trip-sequence}
\end{equation}

The set of all blocks is \(B\), and the vehicle schedule is
\(\Omega=\{b:b\in B\}\). Every selected trip must be covered exactly once:

\begin{equation}
\sum_{b\in B} y_{ib}=1
\qquad
\forall i\in\pi,
\label{eq:problem-trip-coverage}
\end{equation}

where \(y_{ib}=1\) if trip \(i\) is contained in block \(b\). Unselected trips
are not covered:

\begin{equation}
\sum_{b\in B} y_{ib}=0
\qquad
\forall i\in T\setminus\pi.
\label{eq:problem-unselected-trips}
\end{equation}

The number of utilized vehicles equals the number of blocks:

\begin{equation}
N^{veh}=|B|.
\label{eq:problem-vehicle-count}
\end{equation}

The exact-once coverage rule and the exclusion of unselected trips are written
formally in \Cref{eq:trip-covered-once,eq:unselected-not-covered}. These
constraints are the vehicle-schedule side of the timetable--vehicle link.

\subsection{In-Line and Depot-Bridge Compatibility}

Two selected trips \(i\) and \(j\) can be served by the same vehicle only if the
vehicle can finish \(i\) and start \(j\) on time. If \(en(i)=sn(j)\), the
connection is in-line. The stop duration is

\begin{equation}
\Delta_{ij}=st(j)-et(i).
\label{eq:inline-stop-duration-problem}
\end{equation}

The in-line compatibility condition is

\begin{equation}
en(i)=sn(j),
\qquad
\delta^{min}_{en(i),\omega(et(i))}
\le \Delta_{ij}\le
\delta^{max}_{en(i),\omega(et(i))}.
\label{eq:problem-inline-compatibility}
\end{equation}

If the trips do not meet at the same terminal, the vehicle may return to the
depot and leave again. With depot \(dep\), the depot-bridge condition is

\begin{equation}
et(i)
+\tau^{PI}_{en(i),\omega(et(i))}
+\delta^{min}_{dep,\omega(et(i)+\tau^{PI}_{en(i),\omega(et(i))})}
+\tau^{PO}_{sn(j),\omega(st(j))}
\le st(j).
\label{eq:problem-depot-bridge}
\end{equation}

The condition is written with a depot stopping term \(\delta^{min}_{dep}\).
If no positive depot stopping time is specified in an instance,
\(\delta^{min}_{dep}=0\). Keeping it in the equation makes the formulation valid
for instances where such a minimum stop is present. The formal graph edge tests
are given in \Cref{eq:inline,eq:depot}.

The compatibility graph built from these tests is a directed acyclic graph. The
direction comes from time: an edge can only point from an earlier trip to a
later trip. This simplifies vehicle-block construction. A path in the graph is
already a chronological sequence of trips, and a path cover is a set of such
paths that covers the selected timetable. The formal graph definition and its
path-cover use are introduced in \Cref{sec:compatibility-graph-formulation}.

\subsection{Break, Waiting, and Stopping-Time Requirements}

Stopping time is first a feasibility rule. Paid break time is then a cost
component. For a stop activity \(a\), let \(\tau_a\) be the total stop
duration, \(\tau_a^{charge}\) the part used for charging, and
\(\delta^{min}_{n(a),w(a)}\) the required minimum stop at node \(n(a)\) in
the stop-start window \(w(a)\). The paid break part is

\begin{equation}
p_a=
\begin{cases}
\max\left\{0,\tau_a-\max(\tau_a^{charge},\delta^{min}_{n(a),w(a)})\right\},
&\text{in-line stop},\\[1ex]
0,&\text{stop adjacent to pull-in or pull-out}.
\end{cases}
\label{eq:problem-paid-break}
\end{equation}

Thus, the same waiting interval can play different roles. One part may satisfy
minimum stopping time, another part may be used for charging, and only the
remaining part contributes to paid break cost. The corresponding paid-break
constraints are stated in \Cref{eq:paid-break-inline,eq:paid-break-other}.

\section{Electric-Bus and Charging Layer}

\subsection{Battery Autonomy}

For an EV block, residual autonomy is simulated along the ordered activity
sequence. Let \(r_{ba}^{pre}\) and \(r_{ba}^{post}\) be the residual autonomy
before and after activity \(a\) in block \(b\). The vehicle starts from the
depot with full autonomy. For a driving activity,

\begin{equation}
r_{ba}^{post}=r_{ba}^{pre}-\ell_a.
\label{eq:problem-battery-update}
\end{equation}

For a charging activity, \(r_{ba}^{post}=r_{ba}^{pre}+\Delta r_{ba}^{charge}\),
subject to the autonomy upper bound.

An EV assignment is feasible only if

\begin{equation}
0\le r_{ba}^{pre},r_{ba}^{post}\le A^{max}_v
\qquad
\forall a\in A_b.
\label{eq:problem-battery-feasible}
\end{equation}

This explains why EV feasibility is checked after a vehicle block has been
constructed. The same set of trips may be easy or difficult for an EV depending
on the order, deadhead movements, and available charging breaks
\cite{li2014,kooten2017,tang2019}. The formal battery propagation constraints
are given in \Cref{eq:battery-initial,eq:battery-drive,eq:battery-charge,eq:battery-bounds}.

ICE vehicles do not have battery or charging constraints. This does not mean
that an ICE block can ignore the other rules. The selected timetable must first
be feasible, and the block must still satisfy chronological order, pull-in and
pull-out timing, in-line or depot-bridge compatibility, stopping-time rules, and
parking capacity where a break is used. Only after these conventional
vehicle-schedule conditions are satisfied can the battery and charging
restrictions be omitted for an ICE assignment. The EV fleet availability rule is
formalized in \Cref{eq:ev-fleet-limit}; the ICE count relation is given in
\Cref{eq:ice-fleet-count}.

\subsection{Slow and Fast Charging}

Let \(q_a^{slow}\) and \(q_a^{fast}\) be the slow and fast charging durations
inserted during break activity \(a\). Charging must fit inside the available
stop:

\begin{equation}
q_a^{slow}+q_a^{fast}\le \tau_a.
\label{eq:problem-charge-duration}
\end{equation}

For slow charging of vehicle type \(v\), the autonomy gain is

\begin{equation}
\Delta r_a^{slow}
=A^{max}_v\frac{q_a^{slow}}{T^{slow}_v}.
\label{eq:problem-slow-charge}
\end{equation}

For fast charging, the charging time is scaled by the fast-charging multiplier
\(\phi\in(0,1)\):

\begin{equation}
\Delta r_a^{fast}
=A^{max}_v\frac{q_a^{fast}}{\phi_v T^{slow}_v}.
\label{eq:problem-fast-charge}
\end{equation}

The implemented schedule chooses either no charging, slow charging, or fast
charging for a charging activity. This can be expressed as

\begin{equation}
\Delta r_a^{charge}
\in
\left\{0,\Delta r_a^{slow},\Delta r_a^{fast}\right\}.
\label{eq:problem-charge-choice}
\end{equation}

The corresponding duration and gain constraints are written in
\Cref{eq:charge-within-break,eq:charge-upper-slow,eq:charge-upper-fast,eq:charge-gain}.

\subsection{Minimum Charging Time and Partial Recharging}

If a charging activity is used, it must respect the minimum charging duration
of the EV type:

\begin{equation}
q_a^m=0
\quad\text{or}\quad
q_a^m\ge q^{min,m}_v,
\qquad m\in\{\mathrm{slow},\mathrm{fast}\}.
\label{eq:problem-min-charge}
\end{equation}

Partial charging is allowed because the vehicle does not need to return to full
autonomy after every stop. The upper bound is the maximum battery autonomy:

\begin{equation}
r_{ba}^{post}\le A^{max}_v.
\label{eq:problem-charge-cap}
\end{equation}

The minimum charging-time rule and the minute-granularity rule are stated
formally in \Cref{eq:model-min-charge,eq:minute-granularity}.

\subsection{Parking and Charging Capacity at Nodes}

Capacity is checked by overlapping activity intervals. Let \(o_a^{park}\),
\(o_a^{slow}\), and \(o_a^{fast}\) indicate which resource a break window
occupies. Let \(\mathcal{A}_n(t)\) be the break windows active at node \(n\)
and time \(t\). The capacity rules are

\begin{align}
\sum_{a\in\mathcal{A}_n(t)}o_a^{park} &\le K_n^{park}, \label{eq:problem-parking-capacity}\\
\sum_{a\in\mathcal{A}_n(t)}o_a^{slow} &\le K_n^{slow}, \label{eq:problem-slow-capacity}\\
\sum_{a\in\mathcal{A}_n(t)}o_a^{fast} &\le K_n^{fast}. \label{eq:problem-fast-capacity}
\end{align}

These constraints are important because a block can be feasible in isolation
but infeasible together with other blocks if too many vehicles use the same
terminal resource at the same time. The output field \texttt{typeSpot}
specifies the occupied resource, while \texttt{isCharging} specifies whether
active charging occurs in that occupied slot \cite{minoa-io}. The formal
resource-occupation and capacity constraints are given in
\Cref{eq:spot-occupation-vars,eq:active-charging-spot,eq:parking-capacity,eq:slow-capacity,eq:fast-capacity}.

\section{Objective and Cost Structure}

The reported objective is the vehicle-scheduling cost \(C^{VS}\). It contains
fixed vehicle cost, paid break cost, pull-in and pull-out cost, and the
CO\(_2\) cost for conventional driving \cite{minoa-problem}. Fixed vehicle cost is
charged for every used vehicle block. Paid-break cost is charged for residual
break time after mandatory stopping time and active charging have been
accounted for. Pull cost covers the pull-out and pull-in movements between the
depot and the first or last serviced terminal of a block. In the reported
objective decomposition, the CO\(_2\) term applies the vehicle-specific coefficient
to the passenger-service driving duration \(d(b)\). For EV types \(v\in V_E\),
the input coefficient \(c_v^{\mathrm{CO}_2}\) is zero.

\Cref{chap:model} gives the formal objective \(C^{VS}\) in
\Cref{eq:objective} and decomposes it into
\Cref{eq:fixed-cost,eq:break-cost,eq:pull-cost,eq:co2-cost}.

\section{Benchmark Instances and Unit Conventions}

Times are stored in seconds in the input and output files, while result tables
usually report operational durations in minutes:

\begin{equation}
t^{min}=\frac{t^{sec}}{60}.
\label{eq:unit-seconds-minutes}
\end{equation}

All feasibility checks use seconds. Charging windows are written in whole-minute
durations, so a computed required charging duration is rounded up to the next
multiple of 60 seconds before it is inserted into the output. Rounding upward is
used because rounding downward could undercharge an EV block. Capacity is then
checked on these minute-aligned charging and parking intervals. Result tables
divide the final accumulated seconds by 60 for readability; the displayed
minutes do not change the validator-facing schedule.

Distances and autonomy follow the units provided in the instance files. Costs
are reported in the same cost units as the problem statement and the checking
output.

\Cref{tab:unit-conventions} summarizes the unit conventions used in the
formulation, implementation, and reported results.

\begin{table}[htbp]
\centering
\small
\caption{Main unit conventions used in the thesis.}
\label{tab:unit-conventions}
\begin{tabular}{lll}
\toprule
Symbol & Meaning & Unit \\
\midrule
\(\ell_i\) & passenger-trip distance & km \\
\(\ell_a^{dh}\) & deadhead-arc distance & km \\
\(\tau_i^{trip}\) & passenger-trip driving time & s \\
\(\tau_a^{dh}\) & deadhead driving time & s \\
\(d(b)\) & passenger-service driving duration used in the CO\(_2\) audit & s \\
\(A_v^{max}\) & maximum EV autonomy of type \(v\) & km \\
\(q_a\) & charging duration in a break window & s \\
\(c_v^{\mathrm{CO}_2}\) & conventional driving cost coefficient & cost/s \\
\bottomrule
\end{tabular}
\end{table}
\FloatBarrier

The main comparison instances are Small, Medium, and Large. In addition, the
implementation is evaluated on all available Senior instances to check how the
method behaves beyond the headline files. \Cref{tab:instances} lists the twelve
instances used in the thesis. The table reports the structural features that
later influence graph size, vehicle chaining, and EV feasibility: line
directions, terminal nodes, deadhead arcs, headway windows, candidate trips,
planning horizon, EV fleet cap, and EV autonomy.

\begin{figure}[H]
\centering
\includegraphics[width=\textwidth]{figures/fig42_input_solver_object_map.pdf}
\caption{From benchmark data to solver objects.}
\label{fig:input-solver-output}
\end{figure}
\FloatBarrier

\Cref{fig:input-solver-output} is used as the reference map for the following
chapters. The benchmark data are first converted into timetable, graph, and
vehicle-operation objects. The mathematical model formalizes these objects, and
the methodology chapter explains how the solver constructs them.  The main
interpretation is that the solver does not work directly on a flat JSON file.
It first turns the input into structured objects, and each object is responsible
for one part of the MINOA rules.  The raw benchmark properties become solver
objects such as timetable paths, compatibility arcs, vehicle blocks, charging
intervals, and audit outputs.

\begin{table}[htbp]
\centering
\scriptsize
\caption{Structural characteristics of the Senior benchmark instances.}
\label{tab:instances}
\resizebox{\textwidth}{!}{%
\begin{tabular}{lrrrrrrrr}
\toprule
Instance & Dir. & Nodes & Deadhead & Headway & Cand. trips & Horizon & EV cap & Auton. \\
 & & & arcs & windows & & (min) & & (km) \\
\midrule
Small & 2 & 3 & 4 & 6 & 1,064 & 568 & 5 & 90 \\
Medium & 4 & 4 & 6 & 12 & 3,368 & 870 & 5 & 120 \\
Large & 10 & 7 & 12 & 30 & 6,960 & 743 & 5 & 80 \\
Toy Example & 4 & 4 & 6 & 12 & 1,448 & 390 & 5 & 60 \\
1line & 2 & 3 & 4 & 6 & 1,920 & 959 & 5 & 85 \\
1line 6timeWindow & 2 & 3 & 4 & 12 & 1,920 & 959 & 5 & 85 \\
2lines & 4 & 4 & 6 & 12 & 3,840 & 959 & 5 & 85 \\
2lines 6 timeWindows & 4 & 4 & 6 & 24 & 3,840 & 959 & 5 & 85 \\
3lines & 6 & 5 & 8 & 18 & 5,760 & 959 & 5 & 85 \\
3linesTriangle & 6 & 4 & 6 & 18 & 5,760 & 959 & 5 & 85 \\
5lines & 10 & 7 & 12 & 30 & 9,600 & 959 & 5 & 85 \\
8lines & 16 & 10 & 18 & 48 & 15,360 & 959 & 5 & 85 \\
\bottomrule
\end{tabular}
}
\end{table}
\FloatBarrier

The first three rows show why Small, Medium, and Large form a natural headline
set. They increase from 2 to 10 directions and from 1,064 to 6,960 candidate
trips, while EV autonomy also changes. The additional instances extend the
range up to 16 directions and 15,360 candidate trips. They are used to study
robustness and scaling behavior.

The structural columns in \Cref{tab:instances} affect different parts of the
solution method.  The number of directions determines how many independent
timetable-selection sequences have to be built.  The number of terminal nodes
and deadhead arcs determines how many depot-mediated continuations may exist.
The number of headway windows determines how often the service-frequency rule
changes over the planning horizon.  The number of candidate trips controls the
size of the initial timetable-selection problem.

The candidate-trip count is much larger than the number of selected trips in
the final solutions.  This is expected.  The input contains many potential
departures, while the timetable rules require a feasible subset.  For example,
Large contains 6,960 candidate trips, but the final selected timetable contains
260 trips.  The optimization task is therefore not to operate all candidates.
It is to choose a timetable that satisfies the service rules and can be
operated with a feasible vehicle schedule.

The EV fleet cap and autonomy columns explain why the same algorithm can lead
to different EV shares across instances.  A high EV cap does not mean that all
vehicles can become electric.  The vehicle blocks must also fit the autonomy
limit and must contain feasible charging opportunities.  Conversely, an
instance with moderate autonomy can still use several EVs if the constructed
blocks are short enough or if charging breaks appear at useful nodes.

The planning horizon also matters.  A longer horizon does not automatically
mean that the instance is harder.  If trips are spread over a long day, vehicles
may have more opportunities to wait or charge.  If many trips are concentrated
in overlapping time windows, more vehicles are needed even if the horizon is
short.  This is why the computational chapter reports not only total cost but
also selected trips, vehicle count, deadhead minutes, break minutes, and charge
minutes.

The benchmark structure motivates the graph formulation in
\Cref{chap:formulation}.  Each selected trip becomes a node, and every feasible
continuation becomes an arc.  The size and density of this graph depend on the
time windows, terminal layout, and deadhead movements listed in
\Cref{tab:instances}.  The graph is therefore not an abstract modeling device
only; it is a direct representation of the operational structure of each
instance.
